\section{Integrated Discussion}
\label{sec:discussion}

\subsection{RQ1: Model Hierarchy and Question Matching}
\label{subsec:discussion-model-hierarchy-question-matching}

The first research question asked how blood-flow model classes are organized
and what distinguishes them. The answer is that the hierarchy is a set of
retained-state, closure, coupling, and observation choices, not a ranking of
accuracy. Formaggia, Lamponi, and Quarteroni show that section-averaging and
wall simplification define different 1D operators; Shi, Lawford, and Hose
distinguish lumped global pressure--flow interactions from distributed 1D pulse
propagation; and Quarteroni--Veneziani multiscale coupling shows that boundary
data can be a reduced model of omitted circulation
\cite{FormaggiaEtAl2003OneDimensionalBloodFlow,ShiEtAl2011BloodFlowModelReview,QuarteroniVeneziani2003GeometricalMultiscale}.
Thus 0D, 1D, and 2D formulations should not be described as weaker versions of
3D CFD; they retain different states and require different observation operators.

The case study demonstrates this answer by using a distributed 1D
area--flow state only for quantities it can natively or reconstructively
support: area, flow, section-mean velocity, wall-law pressure, and a declared
velocity observation operator. It does not ask that state to produce native WSS,
secondary flow, wall stress, or clinical pressure-ratio evidence.

\subsection{RQ2: Observables and Closure Choices}
\label{subsec:discussion-closures-interpretability}

The second research question asked how rheology, wall mechanics, geometry,
boundary data, numerical discretization, evidence standards, and observation
operators affect interpretation. The answer is that these choices define the
quantity being reported. Blanco and coauthors support this point for
FFR-specific 1D--3D comparison under matched pressure boundary conditions and
dedicated loss closures; Liu and coauthors show observable-dependent rheology
sensitivity in an intracranial stenosis setting; John and coauthors show that
WSS-vector convention changes AWSS/OSI; and Gambaruto and coauthors show that
geometry reconstruction can be a material sensitivity source in a patient
aneurysm case
\cite{Blanco2018FFR1D3D,LiuEtAl2021ICASRheology,JohnEtAl2017WallShearStressVectorForm,GambarutoEtAl2011GeometryRheologySensitivity}.

The case study therefore presents its closure list as mathematical information,
not as implementation detail: Newtonian principal rheology, parabolic
profile closure, an elastic pressure--area wall law, analytic reference
geometry, fixed-area characteristic boundary approximation, and a
plane--tetrahedron velocity observation operator. The resulting comparison is a
velocity-observation comparison under those assumptions, not a WSS,
wall-displacement, or FFR evidence claim.
The membrane discussion sharpens the same boundary: Newtonian or
generalized-Newtonian terminology belongs to the fluid stress law, while the
membrane model is a structural wall closure and the quasi-static comparator is
not evidence for a full transient moving-boundary FSI model.

\subsection{RQ3: Numerical Evidence and Cross-Dimensional Interpretation}
\label{subsec:discussion-numerical-evidence-observation}

The third research question asked what limitations remain when resolved scale,
cost, identifiability, validation status, reproducibility, and uncertainty are
compared. The answer is that interpretation depends on the declared equation,
retained state, discretization, observation operator, and metric, and each
evidence type checks only part of that chain.
ASME VVUQ terminology separates verification, validation, and uncertainty
quantification; Boileau and coauthors provide benchmark comparison evidence;
Pimentel-Garcia and coauthors show that a well-balanced claim must name the
preserved stationary family; and Colebank shows that practical identifiability
depends on measured outputs, boundary-condition class, and design assumptions
\cite{ASME2022VVUQTerminology,BoileauEtAl2015Benchmark,PimentelGarciaEtAl2023HighOrderWellBalanced,Colebank2025HemodynamicsIdentifiability}.

The worked stenosis example demonstrates the distinction. Its manufactured
solution checks the discrete operator for one forced equation. Its
geometry-rest test addresses equilibrium preservation for the chosen
discretization and boundary approximation. Its 1D--3D comparison becomes
meaningful only after both models are reduced to the same section-mean velocity
observable and discrepancy metric. These checks support a bounded diagnostic
cross-model comparison, not an external validation result. The positive MMS
result, negative rest-state result, and final-time comparator are useful
precisely because they belong to different parts of the same interpretive
chain rather than to one validation grade.

\subsection{Contribution and Limitations}
\label{subsec:discussion-contribution-limitations}

The contribution is a review framework for reading mathematical hemodynamics
across model tiers. It connects retained state, closures, boundary data,
discretization, observation operator, metric, and evidence standard in one
interpretive chain, then applies that chain once in an idealized stenosis
computation.

The limitations are correspondingly bounded. The review is narrative and
scoping, not systematic. Some sources remain preprints, software artifacts, or
metadata entries with roles narrower than peer-reviewed empirical evidence. The
case geometry, wall law, rheology, and boundary approximation are idealized.
The numerical evidence distinguishes verification, equilibrium preservation,
and diagnostic comparison, but it does not provide externally validated
clinical prediction evidence or a fully matched transient FSI target.

\section{Conclusion}
\label{sec:conclusions-limitations}

This report's main conclusion is that hemodynamic simulation results are
interpretable only when the equation, retained state, closure assumptions,
boundary data, discretization, observation operator, and metric are declared
together and then evaluated with the appropriate evidence standard. Model
dimension alone does not determine interpretation, and neither a benchmark, a
grid study, a cross-model comparison, nor a learned surrogate result
substitutes for validation against an external target for a specified quantity
and context.
Likewise, a membrane-wall update, a pressure--area law, and a resolved FSI
calculation answer different structural questions unless their wall state,
loads, and observation operators are matched.

The report does not claim externally validated clinical prediction for the
idealized stenosis case or universal superiority of any model tier. Future work
should strengthen the same chain where the present evidence is intentionally
limited: construct or adopt a well-balanced discretization for the
geometry-rest state, obtain matched wall and boundary metadata for resolved
comparisons, report uncertainty separately from model discrepancy, and specify
validation observables before predictive or clinical claims are made.
